\documentclass[12pt,a4paper]{article}

% Encodage et langue
\usepackage[utf8]{inputenc}
\usepackage[T1]{fontenc}
\usepackage[french]{babel}

% Mise en page
\usepackage[margin=2.2cm]{geometry}
\usepackage{fancyhdr}
\usepackage{titlesec}
\usepackage{parskip}

% Mathématiques
\usepackage{amsmath}
\usepackage{amssymb}
\usepackage{bm}

% Code et listings
\usepackage{listings}
\usepackage{xcolor}

% Graphiques et tableaux
\usepackage{graphicx}
\usepackage{booktabs}
\usepackage{array}
\usepackage{float}

% Diagrammes
\usepackage{tikz}
\usetikzlibrary{shapes.geometric,arrows.meta,positioning,calc,fit,backgrounds,patterns}

% Liens
\usepackage{hyperref}

% -----------------------------------------------
% Configuration des listings (code)
% -----------------------------------------------
\definecolor{codegray}{rgb}{0.5,0.5,0.5}
\definecolor{codepurple}{rgb}{0.58,0,0.82}
\definecolor{codeblue}{rgb}{0.0,0.3,0.7}
\definecolor{backcolour}{rgb}{0.97,0.97,0.97}

\lstdefinestyle{cppstyle}{
    backgroundcolor=\color{backcolour},
    commentstyle=\color{codegray}\itshape,
    keywordstyle=\color{codeblue}\bfseries,
    numberstyle=\tiny\color{codegray},
    stringstyle=\color{codepurple},
    basicstyle=\ttfamily\small,
    breakatwhitespace=false,
    breaklines=true,
    captionpos=b,
    keepspaces=true,
    numbers=left,
    numbersep=5pt,
    showspaces=false,
    showstringspaces=false,
    showtabs=false,
    tabsize=2,
    frame=single,
    language=C++
}

\lstset{style=cppstyle}

% -----------------------------------------------
% En-têtes et pieds de page
% -----------------------------------------------
\pagestyle{fancy}
\fancyhf{}
\rhead{CUDA Ray Tracer}
\lhead{Rapport Technique}
\rfoot{Page \thepage}
\lfoot{Mars 2026}

% -----------------------------------------------
% Titre
% -----------------------------------------------
\title{
    \vspace{1cm}
    {\LARGE \textbf{Rapport Technique}} \\[0.5cm]
    {\Large CUDA Ray Tracer} \\[0.3cm]
    {\large Moteur de rendu physiquement réaliste} \\
    {\large GPU (CUDA) / CPU (OpenMP)}
    \vspace{0.5cm}
}
\author{Projet RayTracing}
\date{Mars 2026}

% ===============================================
\begin{document}
% ===============================================

\maketitle
\thispagestyle{empty}

\begin{abstract}
Ce rapport présente l'implémentation d'un moteur de ray tracing physiquement réaliste
tirant parti des capacités de calcul parallèle des GPU modernes via CUDA. Le projet
propose deux chemins de rendu : un chemin GPU utilisant CUDA pour des performances
optimales (~160 MRays/s sur RTX 4060), et un chemin CPU multi-thread utilisant
OpenMP. L'accélération GPU est de l'ordre de \textbf{×56} par rapport à l'exécution CPU.
\end{abstract}

\vspace{0.5cm}
\tableofcontents
\newpage

% ===============================================
\section{Introduction}
% ===============================================

Le ray tracing est une technique de rendu qui simule le comportement physique de la lumière
en traçant des rayons depuis la caméra vers la scène. Contrairement à la rastérisation,
cette approche permet de produire naturellement des effets tels que les réflexions,
les réfractions, les ombres douces et la profondeur de champ.

Ce projet implémente un \textbf{path tracer Monte Carlo} : pour chaque pixel, de nombreux
rayons sont lancés avec des perturbations aléatoires, et leurs contributions sont moyennées
pour obtenir une image convergée. La nature massivement parallèle du problème
(chaque pixel est indépendant) le rend idéal pour les GPU.

\subsection{Objectifs}

\begin{itemize}
    \item Implémenter un path tracer complet avec matériaux physiques (Lambertian, Métal, Diélectrique)
    \item Exploiter CUDA pour une accélération GPU massive
    \item Fournir un chemin de rendu CPU avec OpenMP comme alternative portable
    \item Construire un pipeline de tests unitaires sans dépendance CUDA
\end{itemize}

Le code source complet est disponible sur GitHub :
\url{https://github.com/Gotman08/RayTracing.git}

Le système de build utilise CMake 3.18+ avec trois modes : build GPU complet (défaut),
build CPU seul (\texttt{BUILD\_CPU\_ONLY=ON}), et tests seuls (\texttt{BUILD\_TESTS\_ONLY=ON}).
Les images sont exportées aux formats PNG, BMP, JPG, TGA et PPM.

\subsection{Technologies utilisées}

\begin{table}[H]
\centering
\begin{tabular}{ll}
\toprule
\textbf{Technologie} & \textbf{Rôle} \\
\midrule
CUDA 12.6           & Rendu GPU parallèle \\
OpenMP              & Rendu CPU multi-thread \\
stb\_image\_write   & Export PNG/BMP/JPG/PPM \\
CMake 3.18+         & Système de build \\
CTest               & Tests unitaires \\
\bottomrule
\end{tabular}
\caption{Technologies et bibliothèques du projet}
\end{table}

% ===============================================
\section{Architecture générale}
% ===============================================

Le projet est structuré autour de bibliothèques d'en-têtes (fichiers \texttt{.cuh}) organisées
par domaine fonctionnel, ce qui favorise la réutilisation du code entre les chemins GPU et CPU.

\subsection{Organisation des modules}

\begin{table}[H]
\centering
\begin{tabular}{lp{9cm}}
\toprule
\textbf{Module} & \textbf{Contenu} \\
\midrule
Utilitaires core  & \texttt{vec3}, \texttt{ray}, \texttt{aabb}, \texttt{interval}, \texttt{timer}, générateur aléatoire \\
Géométrie         & \texttt{sphere}, \texttt{plane}, interface \texttt{hittable} \\
Matériaux         & \texttt{lambertian}, \texttt{metal}, \texttt{dielectric} \\
Accélération      & Arbre BVH (variantes GPU et CPU) \\
Rendu             & Renderers GPU/CPU, tone mapping, dispatch de matériaux \\
Caméra            & Caméra perspective avec DoF et motion blur \\
\bottomrule
\end{tabular}
\caption{Organisation modulaire du projet}
\end{table}

\subsection{Trois modes de fonctionnement}

\begin{enumerate}
    \item \textbf{Mode GPU} (défaut) : rendu CUDA, export image
    \item \textbf{Mode CPU} (\texttt{--cpu}) : rendu OpenMP, export image
\end{enumerate}

% ===============================================
\section{Algorithme de path tracing}
% ===============================================

\subsection{Principe du Monte Carlo path tracing}

Pour chaque pixel $(i, j)$, on lance $N$ rayons avec des perturbations aléatoires.
La couleur finale est la moyenne des contributions :

\[
C(i,j) = \frac{1}{N} \sum_{k=1}^{N} \text{trace}(r_k)
\]

Chaque rayon est suivi de rebond en rebond jusqu'à une profondeur maximale $d_{\max}$
ou jusqu'à ce qu'il s'échappe vers l'arrière-plan.

\subsection{Traversée BVH itérative}

L'accélération par \textbf{Bounding Volume Hierarchy} (BVH) réduit la complexité
de l'intersection de $O(n)$ à $O(\log n)$. La traversée est implémentée de façon
itérative avec une pile (profondeur max : 64 nœuds), ce qui est critique pour
les performances GPU où la récursion est coûteuse.

La construction du BVH utilise une stratégie de \textit{median split} sur l'axe
le plus long, avec partitionnement par centroïde.

\subsection{Accumulation et tone mapping}

Les couleurs sont accumulées en virgule flottante (HDR). La compression vers
l'espace SDR utilise le \textbf{tone mapping de Reinhard} :

\[
c' = \frac{c}{1 + c}
\]

Suivi d'une \textbf{correction gamma} $\gamma = 2.2$ :

\[
c_{\text{out}} = c'^{1/\gamma}
\]

% ===============================================
\section{Matériaux}
% ===============================================

\subsection{Lambertian (diffus)}

Le matériau Lambertien modélise les surfaces mattes avec une distribution
cosinus-pondérée :

\[
\bm{d}_{\text{scatter}} = \hat{\bm{n}} + \bm{r}_{\text{unit}}
\]

où $\bm{r}_{\text{unit}}$ est un vecteur unitaire aléatoire. L'atténuation est
la couleur \textit{albedo} du matériau.

\subsection{Métal (réflexion spéculaire)}

La réflexion spéculaire avec flou contrôlé :

\[
\bm{d}_{\text{reflect}} = \bm{v} - 2(\bm{v} \cdot \hat{\bm{n}})\hat{\bm{n}} + f \cdot \bm{r}_{\text{unit}}
\]

où $f \in [0,1]$ est le paramètre de \textit{fuzz} contrôlant le flou des
réflexions (0 = miroir parfait).

\subsection{Diélectrique (verre/réfraction)}

Le diélectrique combine réflexion et réfraction via \textbf{l'approximation de Schlick}
pour les équations de Fresnel :

\[
R(\theta) \approx R_0 + (1 - R_0)(1 - \cos\theta)^5, \quad R_0 = \left(\frac{n_1 - n_2}{n_1 + n_2}\right)^2
\]

La réfraction suit la loi de Snell-Descartes :

\[
n_1 \sin\theta_1 = n_2 \sin\theta_2
\]

La \textbf{réflexion interne totale} est gérée : si $n_1/n_2 \cdot \sin\theta_1 > 1$,
le rayon est entièrement réfléchi. L'indice de réfraction par défaut est $n = 1.5$
(verre).

\begin{table}[H]
\centering
\begin{tabular}{llll}
\toprule
\textbf{Matériau} & \textbf{Modèle} & \textbf{Paramètres} & \textbf{Usage} \\
\midrule
Lambertian  & Diffus cosinus    & albedo (RGB)              & Surfaces mattes \\
Metal       & Réflexion + flou  & albedo (RGB), fuzz [0,1]  & Métaux, miroirs \\
Dielectric  & Snell + Schlick   & IOR (défaut 1.5)          & Verre, eau \\
\bottomrule
\end{tabular}
\caption{Résumé des matériaux implémentés}
\end{table}

% ===============================================
\section{Géométrie}
% ===============================================

\subsection{Sphère}

L'intersection rayon-sphère est résolue analytiquement via l'équation quadratique :

\[
\|\bm{o} + t\bm{d} - \bm{c}\|^2 = r^2
\]

soit $at^2 + bt + c = 0$ avec $a = \|\bm{d}\|^2$,
$b = 2\bm{d} \cdot (\bm{o} - \bm{c})$,
$c = \|\bm{o} - \bm{c}\|^2 - r^2$.

La normale de surface est calculée correctement pour les rayons entrant et sortant
(front/back face). Les coordonnées UV sphériques sont également calculées pour
un potentiel texturage.

La AABB de la sphère est le cube $[\bm{c} - r, \bm{c} + r]$ sur chaque axe.

\subsection{Plan}

Le plan infini est supporté comme primitive géométrique supplémentaire,
permettant la construction de sols et de surfaces plates.

% ===============================================
\section{Caméra}
% ===============================================

La caméra utilise une projection perspective avec FOV vertical configurable (défaut : 20°),
position (\textit{lookfrom}), cible (\textit{lookat}) et vecteur up.

La \textbf{profondeur de champ} est simulée via un modèle de lentille mince :
les rayons sont échantillonnés depuis un disque d'ouverture, pointant vers le plan
de mise au point (défaut : $f/22$, distance 10 unités).
Le \textbf{motion blur} est obtenu en générant les rayons à des instants aléatoires
entre les temps d'ouverture et de fermeture de l'obturateur.

% ===============================================
\section{Implémentation GPU (CUDA)}
% ===============================================

\subsection{Organisation des threads}

Le rendu est organisé en grille 2D de blocs $16 \times 16$ (256 threads/bloc) :

\[
\text{grid} = \left(\left\lceil\frac{W}{16}\right\rceil, \left\lceil\frac{H}{16}\right\rceil\right)
\]

Chaque thread traite un pixel de l'image de manière indépendante.

\begin{figure}[H]
\centering
\begin{tikzpicture}[
    block/.style={draw, minimum width=1.1cm, minimum height=1.1cm, font=\tiny, align=center},
    warp/.style={draw, minimum width=0.28cm, minimum height=0.28cm, font=\tiny},
]
\definecolor{w0}{RGB}{66,133,244}
\definecolor{w1}{RGB}{234,67,53}
\definecolor{w2}{RGB}{251,188,4}
\definecolor{w3}{RGB}{52,168,83}
\definecolor{w4}{RGB}{154,100,215}
\definecolor{w5}{RGB}{0,172,193}
\definecolor{w6}{RGB}{255,112,67}
\definecolor{w7}{RGB}{141,110,99}

% Grid (4x3 blocks)
\node[font=\small\bfseries] at (2, 4.2) {Grid ($\lceil W/16\rceil \times \lceil H/16\rceil$)};
\foreach \bx in {0,1,2,3} {
    \foreach \by in {0,1,2} {
        \pgfmathsetmacro{\highlight}{ifthenelse(\bx==1 && \by==1, 1, 0)}
        \ifnum\bx=1 \ifnum\by=1
            \node[block, fill=green!15, thick, draw=green!50!black] at (\bx*1.3, \by*1.3) {Block\\(\bx,\by)};
        \else
            \node[block, fill=gray!8] at (\bx*1.3, \by*1.3) {(\bx,\by)};
        \fi\else
            \node[block, fill=gray!8] at (\bx*1.3, \by*1.3) {(\bx,\by)};
        \fi
    }
}

% Arrow to zoom
\draw[-{Stealth}, thick, green!50!black] (2.2, 1.3) -- (5.3, 2.5);

% Zoomed block (16x16 = 8 warps)
\node[font=\small\bfseries] at (8.2, 4.2) {Bloc 16$\times$16 = 256 threads};
\node[font=\scriptsize] at (8.2, 3.7) {8 warps de 32 threads};

% Draw 16x16 grid with warp colors (each row of 2 lines = 1 warp)
\foreach \row in {0,...,15} {
    \foreach \col in {0,...,15} {
        \pgfmathtruncatemacro{\warpid}{int((\row*16+\col)/32)}
        \pgfmathtruncatemacro{\cidx}{mod(\warpid, 8)}
        \ifcase\cidx
            \def\wcolor{w0}\or\def\wcolor{w1}\or\def\wcolor{w2}\or\def\wcolor{w3}%
            \or\def\wcolor{w4}\or\def\wcolor{w5}\or\def\wcolor{w6}\or\def\wcolor{w7}\fi
        \fill[\wcolor!35] (5.5+\col*0.19, 3.2-\row*0.19) rectangle ++(0.17, 0.17);
        \draw[gray!30, very thin] (5.5+\col*0.19, 3.2-\row*0.19) rectangle ++(0.17, 0.17);
    }
}

% Warp legend
\foreach \w/\wc in {0/w0, 1/w1, 2/w2, 3/w3, 4/w4, 5/w5, 6/w6, 7/w7} {
    \fill[\wc!45] (11.2, 3.2-\w*0.37) rectangle ++(0.25, 0.25);
    \node[font=\tiny, anchor=west] at (11.55, 3.32-\w*0.37) {warp \w};
}

% Annotation
\node[font=\scriptsize\itshape, text=gray, align=center] at (8.2, -0.3)
    {1 warp = 32 threads contigus = ex\'ecution SIMT (lockstep)};
\end{tikzpicture}
\caption{Organisation hi\'erarchique : Grid $\to$ Blocs $\to$ Warps $\to$ Threads}
\end{figure}

\subsection{Kernels CUDA}

\begin{table}[H]
\centering
\begin{tabular}{lp{8cm}}
\toprule
\textbf{Kernel} & \textbf{Rôle} \\
\midrule
\texttt{init\_rand\_states\_fast}    & Initialisation des états RNG (Wang hash) \\
\texttt{render\_kernel}              & Rendu principal (SPP samples par pixel) \\
\texttt{convert\_to\_rgba8}          & Conversion HDR→RGBA8 avec tone mapping \\
\bottomrule
\end{tabular}
\caption{Kernels CUDA du projet}
\end{table}

\subsection{Pipeline de rendu GPU en 3 passes}

Le rendu GPU s'effectue en trois passes successives sur le framebuffer :

\begin{figure}[H]
\centering
\begin{tikzpicture}[
    node distance=0.4cm,
    kernel/.style={draw, rounded corners, fill=green!15, minimum height=1cm,
                   minimum width=2.2cm, font=\footnotesize\ttfamily, align=center},
    transfer/.style={draw, rounded corners, fill=blue!12, minimum height=1cm,
                     minimum width=2cm, font=\footnotesize\ttfamily, align=center},
    data/.style={font=\scriptsize\itshape, text=gray!70!black},
    arr/.style={-{Stealth[length=2.5mm]}, thick}
]
% Kernels
\node[kernel] (init) {init\_rand\\\_states\_fast};
\node[kernel, right=0.6cm of init] (render) {render\\\_kernel};
\node[kernel, right=0.6cm of render] (reduce) {compute\_avg\\\_luminance};
\node[kernel, right=0.6cm of reduce] (tone) {apply\_tone\\mapping\_kernel};
\node[transfer, right=0.6cm of tone] (d2h) {cudaMemcpy\\Async D$\to$H};
\node[transfer, right=0.5cm of d2h] (png) {save\\PNG};

% Arrows
\draw[arr] (init) -- (render);
\draw[arr] (render) -- node[above, data] {Color* HDR} (reduce);
\draw[arr] (reduce) -- node[above, data] {avg\_lum} (tone);
\draw[arr] (tone) -- node[above, data] {Color* SDR} (d2h);
\draw[arr] (d2h) -- (png);

% Constant memory annotation
\node[below=0.5cm of render, font=\scriptsize, text=purple!70!black, align=center]
    (const) {\texttt{\_\_constant\_\_} Camera + RenderConfig (broadcast cache)};
\draw[-{Stealth}, dashed, purple!50, thin] (const.north) -- (render.south);

% Shared memory annotation
\node[below=0.5cm of reduce, font=\scriptsize, text=orange!70!black, align=center]
    (shared) {\texttt{\_\_shared\_\_} float sdata[256]};
\draw[-{Stealth}, dashed, orange!50, thin] (shared.north) -- (reduce.south);
\end{tikzpicture}
\caption{Pipeline GPU en 3 passes : HDR brut $\to$ r\'eduction luminance $\to$ tone mapping adaptatif}
\end{figure}

\textit{Voir l'animation GIF \texttt{output/anim\_pipeline\_gpu.gif} pour une visualisation pas-\`a-pas.}

\subsection{Gestion mémoire}

\begin{itemize}
    \item \textbf{Host→Device} : matériaux, nœuds BVH, primitives (via \texttt{cudaMemcpy})
    \item \textbf{Device→Host} : framebuffer final après rendu
    \item \textbf{Fixup de pointeurs} : les pointeurs de matériaux host sont redirigés vers
    les adresses device en un seul passage
    \item \textbf{États RNG} : un \texttt{curandState} par pixel (mémoire device)
\end{itemize}

\begin{figure}[H]
\centering
\begin{tikzpicture}[
    box/.style={draw, rounded corners, minimum width=3.8cm, minimum height=0.55cm,
                font=\footnotesize\ttfamily},
    hostbox/.style={box, fill=blue!8},
    devbox/.style={box, fill=green!8},
    arr/.style={-{Stealth[length=2mm]}, thick},
    h2d/.style={arr, blue!70},
    d2h/.style={arr, orange!70},
    lbl/.style={font=\scriptsize, midway},
]

% Host column
\node[font=\bfseries\small] (htitle) {HOST (CPU)};
\node[hostbox, below=0.3cm of htitle] (hmat) {h\_materials[]};
\node[hostbox, below=0.15cm of hmat] (hobj) {h\_objects[]};
\node[hostbox, below=0.15cm of hobj] (hcam) {Camera};
\node[hostbox, below=0.15cm of hcam] (hcfg) {RenderConfig};
\node[hostbox, below=0.5cm of hcfg] (hfb) {h\_frame\_buffer[]};

% Device column
\node[font=\bfseries\small, right=4.5cm of htitle] (dtitle) {DEVICE (GPU)};
\node[devbox, below=0.3cm of dtitle] (dmat) {d\_materials[]};
\node[devbox, below=0.15cm of dmat] (dbvh) {BVH (nodes + prims)};
\node[devbox, below=0.15cm of dbvh, fill=purple!10] (dcam) {\_\_constant\_\_ cam};
\node[devbox, below=0.15cm of dcam, fill=purple!10] (dcfg) {\_\_constant\_\_ cfg};
\node[devbox, below=0.5cm of dcfg] (dfb) {d\_frame\_buffer[]};

% Arrows H→D
\draw[h2d] (hmat.east) -- node[lbl, above] {\scriptsize cudaMemcpy} (dmat.west);
\draw[h2d] (hobj.east) -- node[lbl, above] {\scriptsize BVH build + upload} (dbvh.west);
\draw[h2d] (hcam.east) -- node[lbl, above] {\scriptsize cudaMemcpyToSymbol} (dcam.west);
\draw[h2d] (hcfg.east) -- node[lbl, above] {\scriptsize cudaMemcpyToSymbol} (dcfg.west);

% Arrow D→H
\draw[d2h] (dfb.west) -- node[lbl, below] {\scriptsize cudaMemcpyAsync D$\to$H} (hfb.east);

% Pointer fixup annotation
\node[font=\scriptsize\itshape, text=red!60!black, below=0.05cm of hobj, xshift=2.2cm]
    {pointer fixup H$\to$D};

% Bus annotation
\node[font=\scriptsize, text=gray, below=2.8cm of htitle, xshift=2.3cm]
    {$\longleftrightarrow$ PCI-e bus $\longleftrightarrow$};

% Actions
\node[font=\scriptsize\itshape, below=0.15cm of hfb, text=blue!60!black] {save\_image()};
\node[font=\scriptsize\itshape, below=0.15cm of dfb, text=green!50!black] {render\_kernel()};
\end{tikzpicture}
\caption{Flux de donn\'ees Host $\leftrightarrow$ Device : transferts m\'emoire et types utilis\'es}
\end{figure}

\subsection{Initialisation RNG rapide}

Au lieu de \texttt{curand\_init} (coûteux), le projet utilise un
\textbf{Wang hash} pour une initialisation O(1) par pixel :

\begin{lstlisting}[caption={Wang hash pour l'initialisation RNG}]
__device__ uint32_t wang_hash(uint32_t seed) {
    seed = (seed ^ 61) ^ (seed >> 16);
    seed *= 9;
    seed = seed ^ (seed >> 4);
    seed *= 0x27d4eb2d;
    seed = seed ^ (seed >> 15);
    return seed;
}
\end{lstlisting}

% ===============================================
\section{Implémentation CPU (OpenMP)}
% ===============================================

\subsection{Parallélisation}

Le rendu CPU utilise OpenMP avec un ordonnancement dynamique par lignes :

\begin{lstlisting}[caption={Parallélisation OpenMP}]
#pragma omp parallel for schedule(dynamic, 1)
for (int j = 0; j < height; j++) {
    CPURandom rng(seed + j);
    for (int i = 0; i < width; i++) {
        // rendu du pixel (i, j)
    }
}
\end{lstlisting}

L'\textbf{ordonnancement dynamique} (\texttt{schedule(dynamic, 1)}) assure un
équilibrage de charge optimal, certaines lignes étant plus coûteuses que d'autres
selon la complexité de la scène.

\begin{figure}[H]
\centering
\begin{tikzpicture}[
    cell/.style={draw, minimum width=0.85cm, minimum height=0.55cm, font=\scriptsize},
    lbl/.style={font=\footnotesize\bfseries, anchor=east},
]
\definecolor{t0}{RGB}{66,133,244}
\definecolor{t1}{RGB}{234,67,53}
\definecolor{t2}{RGB}{251,188,4}
\definecolor{t3}{RGB}{52,168,83}

% Title
\node[font=\small\bfseries] at (4.5, 1.2) {Ordonnancement dynamique --- schedule(dynamic, 1)};

% Thread labels
\node[lbl] at (-0.1, 0) {Thread 0};
\node[lbl] at (-0.1, -0.7) {Thread 1};
\node[lbl] at (-0.1, -1.4) {Thread 2};
\node[lbl] at (-0.1, -2.1) {Thread 3};

% Thread 0 - lines
\foreach \x/\line in {0/0, 1/5, 2/8, 3/12, 4/15} {
    \node[cell, fill=t0!30] at (0.5+\x*0.95, 0) {\line};
}
% Thread 1
\foreach \x/\line in {0/1, 1/4, 2/9, 3/11, 4/16} {
    \node[cell, fill=t1!30] at (0.5+\x*0.95, -0.7) {\line};
}
% Thread 2
\foreach \x/\line in {0/2, 1/6, 2/7, 3/13, 4/17} {
    \node[cell, fill=t2!30] at (0.5+\x*0.95, -1.4) {\line};
}
% Thread 3
\foreach \x/\line in {0/3, 1/10, 2/14, 3/18, 4/19} {
    \node[cell, fill=t3!30] at (0.5+\x*0.95, -2.1) {\line};
}

% Arrow time
\draw[-{Stealth}, thick, gray] (5.5, 0.5) -- (8, 0.5) node[midway, above, font=\scriptsize] {temps};

% Post-processing
\node[draw, rounded corners, fill=gray!10, minimum width=2.5cm, minimum height=0.6cm,
      font=\scriptsize] at (7, -1.05) {Tone mapping + Gamma};
\draw[-{Stealth}, thick] (5.3, -1.05) -- (5.7, -1.05);

% Note
\node[font=\scriptsize\itshape, text=gray] at (4.5, -3) {Chaque cellule = 1 ligne de l'image, assign\'ee dynamiquement au prochain thread libre};
\end{tikzpicture}
\caption{Pipeline CPU : ordonnancement dynamique OpenMP par lignes}
\end{figure}

\subsection{Générateur aléatoire CPU}

Un générateur à congruence linéaire (\texttt{CPURandom}) est utilisé,
initialisé par ligne pour garantir la reproductibilité et éviter les
conditions de course entre threads.

\subsection{Comparaison GPU vs CPU}

\begin{table}[H]
\centering
\begin{tabular}{lrr}
\toprule
\textbf{Paramètre} & \textbf{GPU (RTX 4060)} & \textbf{CPU (8 threads)} \\
\midrule
Résolution          & 800×600     & 800×600 \\
Samples per pixel   & 100         & 100 \\
Temps de rendu      & $\sim$0.3 s & $\sim$17 s \\
Débit               & $\sim$160 MRays/s & $\sim$2.8 MRays/s \\
\textbf{Speedup GPU/CPU} & \multicolumn{2}{c}{\textbf{×56}} \\
\bottomrule
\end{tabular}
\caption{Comparaison des performances GPU vs CPU}
\end{table}

% ===============================================
\section{Tests unitaires}
% ===============================================

La suite de tests est compilable \textbf{sans CUDA} (\texttt{BUILD\_TESTS\_ONLY=ON}),
ce qui permet de valider les composants mathématiques et géométriques de manière
indépendante du GPU.

\begin{table}[H]
\centering
\begin{tabular}{ll}
\toprule
\textbf{Fichier} & \textbf{Composant testé} \\
\midrule
\texttt{test\_vec3.cpp}     & Opérations vectorielles (construction, arithmétique, normalisation, produits) \\
\texttt{test\_ray.cpp}      & Paramétrisation des rayons \\
\texttt{test\_interval.cpp} & Intersection et containment d'intervalles \\
\texttt{test\_aabb.cpp}     & Intersection et fusion de boîtes englobantes \\
\texttt{test\_sphere.cpp}   & Intersection rayon-sphère \\
\texttt{test\_plane.cpp}    & Géométrie des plans \\
\texttt{test\_materials.cpp} & Matériaux (Lambertian, Métal, Diélectrique) \\
\texttt{test\_camera.cpp}    & Génération de rayons, FOV, profondeur de champ \\
\texttt{test\_tone\_mapping.cpp} & Tone mapping Reinhard et correction gamma \\
\texttt{test\_bvh.cpp}       & Construction et traversée du BVH \\
\texttt{test\_main.cpp}      & Framework de tests et statistiques \\
\bottomrule
\end{tabular}
\caption{Suite de tests unitaires}
\end{table}

La compilation et l'exécution se font via \texttt{cmake .. -DBUILD\_TESTS\_ONLY=ON \&\& make \&\& ctest}.
L'ensemble de la suite CPU comprend \textbf{144 tests} validant toutes les classes du projet.

\subsection{Tests GPU (CUDA)}

Une suite de tests GPU séparée (\texttt{test\_cuda\_kernels.cu}) valide les kernels
spécifiques au GPU. Elle nécessite un GPU NVIDIA et n'est compilée que lorsque
CUDA est disponible :

\begin{table}[H]
\centering
\begin{tabular}{lp{8cm}}
\toprule
\textbf{Test} & \textbf{Vérification} \\
\midrule
\texttt{init\_rand\_states}            & Initialisation curand sans crash \\
\texttt{luminance\_reduction\_red}     & Réduction sur pixels rouges $\rightarrow$ 0.2126 \\
\texttt{luminance\_reduction\_white}   & Réduction sur pixels blancs $\rightarrow$ 1.0 \\
\texttt{reduction\_non\_aligned}       & Buffer non multiple de 256 threads \\
\texttt{adaptive\_tonemapping}         & Sortie dans $[0, 1]$ pour tout input HDR \\
\texttt{tonemapping\_black}            & Noir reste noir après tone mapping \\
\texttt{constant\_memory\_camera}      & Lecture Camera depuis \texttt{\_\_constant\_\_} \\
\texttt{constant\_memory\_config}      & Lecture RenderConfig depuis \texttt{\_\_constant\_\_} \\
\bottomrule
\end{tabular}
\caption{Tests unitaires GPU}
\end{table}

% ===============================================
\section{Optimisations de performance}
% ===============================================

\begin{enumerate}
    \item \textbf{Wang hash RNG} : initialisation O(1) des états aléatoires par pixel,
    évitant le coûteux \texttt{curand\_init}

    \item \textbf{Traversée BVH itérative} : pile explicite (64 nœuds max) évitant
    la récursion, favorable au cache GPU et au parallélisme SIMT

    \item \textbf{Mémoire constante} : Camera et RenderConfig en \texttt{\_\_constant\_\_}
    pour le broadcast cache ($\sim$5 cycles vs $\sim$500 en globale)

    \item \textbf{Tone mapping adaptatif en 2 passes} : réduction parallèle de la
    luminance moyenne (\texttt{\_\_shared\_\_} + \texttt{atomicAdd}), puis
    Reinhard étendu avec exposition automatique

    \item \textbf{CUDA Streams} : recouvrement de l'upload matériaux et de l'init
    RNG sur deux streams indépendants

    \item \textbf{Fast math} : flags \texttt{-{}-use\_fast\_math} (GPU) et \texttt{-ffast-math}
    (CPU) pour des opérations FPU approximatives mais plus rapides

    \item \textbf{Intersection AABB sans branchement} : code déroulé pour minimiser
    les divergences de warp sur GPU

    \item \textbf{Ordonnancement dynamique CPU} : \texttt{schedule(dynamic, 1)} pour
    un équilibrage de charge optimal entre threads OpenMP
\end{enumerate}

% ===============================================
\section{Analyse comparative CPU vs GPU}
% ===============================================

\subsection{Intérêts de la solution GPU (CUDA)}

Le GPU est idéal pour le path tracing grâce à son architecture massivement parallèle :
des milliers de cœurs CUDA traitent chaque pixel de manière indépendante.
La RTX 4060 (3072 cœurs CUDA) atteint \textbf{160 MRays/s}, soit un débit
56 fois supérieur au CPU.

\textbf{Avantages} : débit très élevé, adapté aux problèmes embarrassingly parallel,
mémoire GPU à haute bande passante (GDDR6).

\textbf{Inconvénients} : transferts mémoire Host$\leftrightarrow$Device coûteux, gestion mémoire
complexe (fixup de pointeurs, allocation explicite), débogage plus difficile,
portabilité limitée aux GPU NVIDIA.

\subsection{Intérêts de la solution CPU (OpenMP)}

Le CPU offre portabilité et simplicité. OpenMP permet une parallélisation
à faible coût de développement (une seule directive pragma), et le débogage
se fait avec des outils standards (GDB, Valgrind).

\textbf{Avantages} : portabilité universelle, mémoire partagée (pas de transferts),
débogage direct, support automatique de tous les compilateurs C++ modernes.

\textbf{Inconvénients} : parallélisme limité à 8--16 threads (vs milliers sur GPU),
débit de seulement \textbf{2.8 MRays/s} sur 8 threads.

\subsection{Discussion du speedup effectif}

Le speedup mesuré de \textbf{×56} s'explique par la nature embarrasingly parallel
du path tracing : chaque pixel étant totalement indépendant, le GPU exploite pleinement
ses milliers de cœurs sans synchronisation. Ce ratio est cohérent avec la littérature
sur les accélérations GPU pour les applications de rendu Monte Carlo.

Le mode \texttt{-{}-benchmark} permet de vérifier cette tendance sur plusieurs résolutions,
confirmant le caractère compute-bound du kernel de rendu.

% ===============================================
\section{Hiérarchie mémoire CUDA}
% ===============================================

L'architecture CUDA expose plusieurs niveaux de mémoire, chacun avec des caractéristiques
de latence, de taille et de portée différentes. Le choix du type de mémoire a un impact
direct sur les performances.

\subsection{Vue d'ensemble des types de mémoire}

\begin{table}[H]
\centering
\begin{tabular}{lrrlp{4.5cm}}
\toprule
\textbf{Type}      & \textbf{Taille}   & \textbf{Latence}    & \textbf{Portée}    & \textbf{Usage typique} \\
\midrule
Registres          & $\sim$256 KB/SM   & $\sim$0 cycles      & Thread             & Variables locales \\
Shared (\texttt{\_\_shared\_\_}) & 48--100 KB/SM & $\sim$5 cycles & Bloc & Réductions, stencils \\
Constant (\texttt{\_\_constant\_\_}) & 64 KB total & $\sim$5 cycles (cache) & Grille & Params read-only, broadcast \\
Texture/L1         & $\sim$128 KB/SM   & $\sim$5 cycles (cache) & Grille          & Accès 2D spatialement locaux \\
Globale            & 8--80 GB          & $\sim$400--500 cycles & Grille           & Framebuffers, BVH, données \\
\bottomrule
\end{tabular}
\caption{Caractéristiques des différents types de mémoire CUDA}
\end{table}

\begin{figure}[H]
\centering
\begin{tikzpicture}[
    layer/.style={draw, rounded corners=3pt, minimum height=0.8cm, align=center, font=\footnotesize},
    annot/.style={font=\scriptsize\itshape, text=gray!70!black, anchor=west},
]

% Layers (pyramid: narrower at top, wider at bottom)
\node[layer, fill=green!35, minimum width=4cm] (reg) at (0, 4) {Registres};
\node[layer, fill=green!20, minimum width=5.5cm] (sh)  at (0, 3) {\texttt{\_\_shared\_\_}};
\node[layer, fill=yellow!25, minimum width=7cm] (cst) at (0, 2) {\texttt{\_\_constant\_\_}};
\node[layer, fill=orange!15, minimum width=8.5cm] (l1)  at (0, 1) {L1 / Texture cache};
\node[layer, fill=red!12, minimum width=10cm] (glb) at (0, 0) {M\'emoire globale (VRAM)};

% Left: speed labels
\node[font=\scriptsize, text=green!50!black, anchor=east] at (-5.3, 4) {$\sim$0 cycles};
\node[font=\scriptsize, text=green!50!black, anchor=east] at (-5.3, 3) {$\sim$5 cycles};
\node[font=\scriptsize, text=yellow!60!black, anchor=east] at (-5.3, 2) {$\sim$5 cycles};
\node[font=\scriptsize, text=orange!60!black, anchor=east] at (-5.3, 1) {$\sim$5 cycles};
\node[font=\scriptsize, text=red!50!black, anchor=east] at (-5.3, 0) {400+ cycles};

% Right: usage annotations
\node[annot] at (2.3, 4) {ray, rec, attenuation, BVH stack[64]};
\node[annot] at (3, 3) {sdata[256] (r\'eduction luminance)};
\node[annot] at (3.7, 2) {Camera (120 B), RenderConfig (64 B)};
\node[annot] at (4.5, 1) {cache automatique (acc\`es BVH)};
\node[annot] at (5.3, 0) {framebuffer, BVH, curandState, mat\'eriaux};

% Side labels
\draw[{Stealth}-{Stealth}, thick, gray] (-5.8, -0.3) -- (-5.8, 4.3);
\node[font=\scriptsize, text=gray, rotate=90, anchor=center] at (-6.3, 2) {Rapide $\leftarrow$ Latence $\rightarrow$ Lent};

\draw[{Stealth}-{Stealth}, thick, gray] (10.3, -0.3) -- (10.3, 4.3);
\node[font=\scriptsize, text=gray, rotate=90, anchor=center] at (10.8, 2) {Petit $\leftarrow$ Taille $\rightarrow$ Grand};
\end{tikzpicture}
\caption{Pyramide m\'emoire CUDA : compromis latence / capacit\'e avec les donn\'ees du projet}
\end{figure}

\subsection{Mémoire constante : Camera et RenderConfig}

La \texttt{Camera} ($\sim$120 octets) et la \texttt{RenderConfig} ($\sim$64 octets) sont
lues par \textbf{tous les threads de manière identique}. C'est le cas d'usage idéal
de la mémoire constante : le cache constant diffuse la valeur à tout le warp
en un seul cycle (broadcast), alors qu'un accès en mémoire globale coûte 400--500 cycles.

\begin{lstlisting}[caption={Déclaration et utilisation de la mémoire constante}]
// declaration dans renderer.cuh (namespace scope)
__constant__ Camera d_const_camera;
__constant__ RenderConfig d_const_config;

// copie avant le lancement du kernel
cudaMemcpyToSymbol(d_const_camera, &camera,
                   sizeof(Camera));
\end{lstlisting}

Taille totale : $\sim$184 octets, bien en dessous de la limite de 64 KB.

\subsection{Mémoire partagée : réduction de luminance}

La mémoire partagée (\texttt{\_\_shared\_\_}) est utilisée dans le kernel de réduction
parallèle pour calculer la luminance moyenne. Chaque bloc de 256 threads utilise
un tableau partagé de 1~KB ($256 \times 4$ octets) pour effectuer une réduction
en arbre binaire avec \texttt{\_\_syncthreads()}.

\subsection{Mémoire de texture : discussion}

Les nœuds BVH pourraient théoriquement bénéficier de la mémoire de texture
grâce à son cache optimisé pour les accès spatialement locaux. Cependant,
les architectures récentes (sm\_70+) unifient les caches L1 et texture,
réduisant l'avantage. De plus, la traversée BVH suit un ordre dépendant
des rayons (accès irrégulier), ce qui limite l'efficacité du cache de texture.
Le texturage serait plus pertinent pour un ray tracer avec textures UV mappées.

\subsection{Coalescence des accès mémoire}

La \textbf{coalescence} est la propriété par laquelle les threads d'un même warp accèdent
à des adresses mémoire consécutives, permettant de fusionner les accès en une seule
transaction mémoire large (128 octets sur les GPU modernes).

Dans notre implémentation, l'index linéaire du pixel est calculé ainsi :
\[
\texttt{pixel\_index} = j \times W + i
\]
Avec les threads organisés en $x$ (dimension rapide du bloc $16\times16$), les threads
$t_0, t_1, \ldots, t_{15}$ d'un warp correspondent à des colonnes consécutives
($i, i+1, \ldots, i+15$) sur la même ligne $j$. Leurs adresses dans le framebuffer
et le tableau \texttt{curandState} sont donc \textbf{contiguës} → accès coalescés,
bande passante maximale.

À l'inverse, un accès par colonnes ($j$ comme dimension rapide) produirait des
accès avec stride $W$, cassant la coalescence et réduisant le débit d'un facteur
pouvant dépasser $\times$10.

% ===============================================
\section{Réduction parallèle et Thrust}
% ===============================================

\subsection{Algorithme de réduction en arbre binaire}

Pour calculer la luminance moyenne de l'image (nécessaire au tone mapping adaptatif),
on implémente une \textbf{réduction parallèle}. L'algorithme procède en trois étapes :

\begin{enumerate}
    \item Chaque thread charge la luminance de son pixel en \texttt{\_\_shared\_\_} memory
    \item Réduction intra-bloc par divisions successives avec \texttt{\_\_syncthreads()}
    \item Le thread 0 de chaque bloc accumule son résultat via \texttt{atomicAdd}
\end{enumerate}

\begin{lstlisting}[caption={Kernel de réduction parallèle (extrait)}]
__global__ void compute_avg_luminance(
    const Color* frame_buffer,
    float* d_total_luminance, int num_pixels
) {
    __shared__ float sdata[256];
    int tid = threadIdx.x;
    int gid = blockIdx.x * blockDim.x + threadIdx.x;

    sdata[tid] = (gid < num_pixels) ?
        luminance(frame_buffer[gid]) : 0.0f;
    __syncthreads();

    for (unsigned int s = blockDim.x / 2; s > 0; s >>= 1) {
        if (tid < s) sdata[tid] += sdata[tid + s];
        __syncthreads();
    }

    if (tid == 0) atomicAdd(d_total_luminance, sdata[0]);
}
\end{lstlisting}

Complexité : $O(N/B + \log B)$ par bloc, où $N$ = nombre de pixels et $B$ = 256.

\subsection{Comparaison avec Thrust}

La librairie Thrust (fournie avec le CUDA Toolkit) propose
\texttt{thrust::transform\_reduce} qui effectue la même opération en une seule ligne :

\begin{lstlisting}[caption={Réduction avec Thrust}]
float total = thrust::transform_reduce(
    d_ptr, d_ptr + num_pixels,
    LuminanceFunctor(), 0.0f, thrust::plus<float>());
\end{lstlisting}

\textbf{Avantage du kernel custom} : contrôle fin sur la shared memory et registres,
optimisations spécifiques (\texttt{\_\_shfl\_down\_sync}, loop unrolling).

\textbf{Avantage de Thrust} : code plus concis, optimisé par NVIDIA, portable
entre architectures.

\begin{figure}[H]
\centering
\includegraphics[width=0.65\textwidth]{../output/chart_reduction_comparison.png}
\caption{Comparaison des temps de réduction : kernel custom (warp shuffle) vs Thrust}
\end{figure}

\subsection{Optimisation warp-level : \texttt{\_\_shfl\_down\_sync}}

La réduction en arbre binaire nécessite un \texttt{\_\_syncthreads()} à chaque itération.
Or, les 5 dernières itérations (strides 32, 16, 8, 4, 2, 1) impliquent uniquement
les threads 0--31, qui forment un seul warp. Dans un warp, tous les threads
s'exécutent en lockstep : \texttt{\_\_syncthreads()} y est redondant et coûteux.

L'instruction \texttt{\_\_shfl\_down\_sync} permet d'échanger des valeurs de \textbf{registres}
directement entre threads du même warp, sans passer par la shared memory :

\begin{lstlisting}[caption={Réduction warp-level avec \_\_shfl\_down\_sync}]
if (tid < 32) {
    float val = sdata[tid] + sdata[tid + 32]; // stride=32
    val += __shfl_down_sync(0xffffffff, val, 16);
    val += __shfl_down_sync(0xffffffff, val, 8);
    val += __shfl_down_sync(0xffffffff, val, 4);
    val += __shfl_down_sync(0xffffffff, val, 2);
    val += __shfl_down_sync(0xffffffff, val, 1);
    if (tid == 0) atomicAdd(d_total_luminance, val);
}
\end{lstlisting}

Le masque \texttt{0xffffffff} signifie que les 32 threads du warp participent tous.
Cette optimisation supprime 5 \texttt{\_\_syncthreads()} et 5 séries d'accès shared memory,
réduisant la latence de la réduction intra-bloc.

\textit{Voir l'animation GIF \texttt{output/anim\_reduction.gif} pour une visualisation stride par stride.}

\subsection{Note sur cuBLAS}

cuBLAS est la librairie CUDA pour l'algèbre linéaire dense (BLAS Level 1/2/3).
Dans un ray tracer, son utilité directe est limitée car les opérations sont
per-ray (pas de grandes matrices). Elle pourrait être pertinente pour des
transformations de géométrie instanciée (matrices $4\times4$ par batch d'objets).

% ===============================================
\section{CUDA Streams et recouvrement}
% ===============================================

\subsection{Principe des streams}

Un \texttt{cudaStream\_t} représente une file d'exécution asynchrone sur le GPU.
Les opérations dans des streams différents peuvent s'exécuter en parallèle,
permettant de recouvrir les transferts mémoire avec le calcul.

\subsection{Implémentation dans le projet}

Nous utilisons deux streams :

\begin{itemize}
    \item \texttt{stream\_compute} : kernels de calcul (init RNG, rendu)
    \item \texttt{stream\_transfer} : transferts mémoire (H2D matériaux, D2H framebuffer)
\end{itemize}

\textbf{Recouvrement 1} : l'upload des matériaux (\texttt{cudaMemcpyAsync} sur
\texttt{stream\_transfer}) s'exécute en parallèle de l'initialisation des états
curand (\texttt{init\_rand\_states\_fast} sur \texttt{stream\_compute}).
Les deux opérations sont indépendantes et opèrent sur des zones mémoire différentes.

\textbf{Recouvrement 2} : après le rendu, le téléchargement du framebuffer est
lancé en asynchrone pendant que le CPU prépare la sauvegarde du fichier.

Le gain est faible car le kernel de rendu domine largement le temps total,
mais cette technique démontre la maîtrise du concept d'exécution asynchrone.

\begin{figure}[H]
\centering
\begin{tikzpicture}[
    block/.style={draw, rounded corners=2pt, minimum height=0.6cm, font=\scriptsize\ttfamily, align=center},
    comp/.style={block, fill=green!20},
    xfer/.style={block, fill=blue!15},
    sync/.style={draw, dashed, red!60, thick},
    overlap/.style={draw=none, fill=orange!10, rounded corners=2pt},
    lbl/.style={font=\footnotesize\bfseries, anchor=east},
]

% Labels
\node[lbl] at (-0.2, 0) {stream\_compute};
\node[lbl] at (-0.2, -1.2) {stream\_transfer};

% Time axis
\draw[-{Stealth}, gray, thick] (0, -2.2) -- (14, -2.2) node[right, font=\scriptsize] {temps};
\foreach \x in {0,2,...,12} { \draw[gray, thin] (\x, -2.1) -- (\x, -2.3); }

% Overlap zones
\fill[orange!8] (0, -1.55) rectangle (2.5, 0.35);
\node[font=\tiny\itshape, text=orange!70!black] at (1.25, -1.85) {overlap 1};

\fill[orange!8] (12, -1.55) rectangle (14, 0.35);
\node[font=\tiny\itshape, text=orange!70!black] at (13, -1.85) {overlap 2};

% stream_compute
\node[comp, minimum width=2cm] at (1, 0) {init\_rand\\\_states};
\node[comp, minimum width=6cm] at (6.5, 0) {render\_kernel};
\node[comp, minimum width=1.3cm] at (10.5, 0) {reduce};
\node[comp, minimum width=1.3cm] at (12, 0) {tonemap};

% stream_transfer
\node[xfer, minimum width=2cm] at (1, -1.2) {materials\\H$\to$D};
\node[xfer, minimum width=1.5cm] at (3, -1.2) {BVH\\H$\to$D};
\node[xfer, minimum width=1.5cm] at (13.2, -1.2) {FB\\D$\to$H};

% Sync points
\draw[sync] (2.5, 0.4) -- (2.5, -1.6) node[below, font=\tiny, text=red!60] {sync};
\draw[sync] (9.5, 0.4) -- (9.5, -0.4);
\end{tikzpicture}
\caption{Diagramme temporel des 2 CUDA streams : zones de recouvrement en orange}
\end{figure}

\textit{Voir l'animation GIF \texttt{output/anim\_streams.gif} pour la progression temporelle des streams.}

% ===============================================
\section{Contrôle de flux et divergence de warp}
% ===============================================

Sur un GPU NVIDIA, les threads sont exécutés par groupes de 32 (un \textbf{warp}).
Lorsque les threads d'un même warp prennent des branches conditionnelles différentes,
on parle de \textbf{divergence de warp}. Les deux branches sont exécutées
séquentiellement, divisant l'efficacité.

\subsection{Sources de divergence identifiées}

\begin{enumerate}
    \item \textbf{Dispatch de matériaux} : le \texttt{switch(mat.type)} dans
    \texttt{scatter()} cause de la divergence quand des pixels voisins d'un même
    warp touchent des matériaux différents (ex : un pixel frappe une sphère métallique,
    son voisin un sol lambertien). La divergence est maximale aux frontières
    entre objets de matériaux différents.

    \item \textbf{Traversée BVH} : le test \texttt{if (node.is\_leaf)} provoque
    de la divergence car les rayons d'un même warp peuvent atteindre des profondeurs
    différentes dans l'arbre. De plus, certains rayons s'échappent vers le ciel
    après 2 rebonds, d'autres en font 20.

    \item \textbf{Guard de bordure} : \texttt{if (i >= width || j >= height) return}
    n'affecte que les threads de bord (impact négligeable).
\end{enumerate}

\subsection{Atténuation}

L'utilisation d'un dispatch par \texttt{enum/switch} au lieu de fonctions virtuelles
permet au compilateur de prédire les branches. Un tri des rayons par matériau
(material sorting) pourrait réduire la divergence du cas 1, mais au prix
d'une synchronisation globale coûteuse. Pour notre scène, le ratio performance/complexité
ne justifie pas cette optimisation.

\textit{Voir l'animation GIF \texttt{output/anim\_warp\_divergence.gif} pour une visualisation de l'impact du switch(mat.type).}

% ===============================================
\section{Analyse d'occupancy et profilage GPU}
% ===============================================

\subsection{Occupancy théorique}

L'\textbf{occupancy} mesure le ratio entre les warps actifs et le maximum de warps
supportés par SM. Elle est calculée via l'API :

\begin{lstlisting}[caption={Calcul d'occupancy via l'API CUDA}]
int num_blocks_per_sm;
cudaOccupancyMaxActiveBlocksPerMultiprocessor(
    &num_blocks_per_sm, render_kernel,
    BLOCK_SIZE * BLOCK_SIZE, 0);
// occupancy = (active_warps / max_warps_per_sm)
\end{lstlisting}

Avec des blocs de $16 \times 16 = 256$ threads (8 warps/bloc), l'occupancy
dépend du nombre de registres utilisés par le kernel et de la shared memory consommée.
L'API \texttt{cudaOccupancyMaxPotentialBlockSize} suggère la taille de bloc optimale.

\subsection{Estimation de la bande passante}

La bande passante effective est estimée par :
\[
BW = \frac{\text{bytes lus + écrits}}{\text{temps du kernel}}
\]

Pour le kernel de rendu : chaque pixel lit un \texttt{curandState} ($\sim$52 octets)
et écrit un \texttt{Color} (12 octets), soit $\sim$64 octets/pixel.
Avec les accès BVH et matériaux, la bande passante effective mesurée reste
bien en dessous du pic théorique (ex : 272 GB/s pour une RTX 4060),
ce qui indique que le kernel est \textbf{compute-bound} (limité par le calcul,
pas par la mémoire).

\subsection{Mode --profile}

Le flag \texttt{-{}-profile} affiche un rapport complet incluant :
occupancy par kernel, bande passante estimée, comparaison des temps
de réduction (custom vs Thrust), et tailles des structures en mémoire constante.

\subsection{Mode --benchmark : mesures multi-résolution}

Le flag \texttt{-{}-benchmark} lance automatiquement le rendu sur 4 résolutions
avec 50 SPP chacune, mesuré avec des CUDA events (précision microseconde).
Il permet d'observer la montée en charge du GPU et la saturation du débit.

\begin{table}[H]
\centering
\begin{tabular}{lrrr}
\toprule
\textbf{Résolution} & \textbf{Pixels} & \textbf{Temps (ms)} & \textbf{MRays/s} \\
\midrule
640$\times$360    & 230K   & 3.5    & 705  \\
1280$\times$720   & 922K   & 7.5    & 1265 \\
1920$\times$1080  & 2.1M   & 15.3   & 1399 \\
2560$\times$1440  & 3.7M   & 26.8   & 1443 \\
\bottomrule
\end{tabular}
\caption{Résultats benchmark multi-résolution sur NVIDIA GH200 H100 (50 SPP)}
\end{table}

La saturation à haute résolution ($\sim$1400--1450 MRays/s) confirme que le GPU
est \textbf{compute-bound} : doubler le nombre de pixels double presque le temps
de calcul, le débit restant constant.

\begin{figure}[H]
\centering
\includegraphics[width=0.75\textwidth]{../output/chart_benchmark.png}
\caption{Throughput (MRays/s) en fonction de la résolution - saturation du GPU visible}
\end{figure}

% ===============================================
\section{Conclusion}
% ===============================================

Ce projet démontre une implémentation complète et optimisée d'un path tracer
physiquement réaliste exploitant le parallélisme massif des GPU modernes.

Les points forts du projet sont :

\begin{itemize}
    \item \textbf{Double implémentation GPU/CPU} offrant flexibilité et portabilité
    \item \textbf{Accélération ×56} grâce à CUDA par rapport au rendu CPU
    \item \textbf{Exploitation de la hiérarchie mémoire} : \texttt{\_\_constant\_\_} pour Camera/Config,
    \texttt{\_\_shared\_\_} pour la réduction, mémoire globale pour les données volumineuses
    \item \textbf{Réduction parallèle hybride} : arbre binaire + \texttt{\_\_shfl\_down\_sync}
    (warp-level), comparée à Thrust \texttt{transform\_reduce}
    \item \textbf{Tone mapping adaptatif} (Reinhard étendu) exploitant la luminance moyenne
    \item \textbf{CUDA Streams} pour le recouvrement calcul/transfert
    \item \textbf{Analyse d'occupancy} via l'API CUDA avec estimation de bande passante
    et mode \texttt{-{}-benchmark} multi-résolution (4 résolutions, CUDA events)
    \item \textbf{Coalescence mémoire} assurée par l'indexation $j \times W + i$ (accès contigus)
    \item \textbf{Tests CPU et GPU} : 144 tests CPU + 8 tests de kernels CUDA
    \item \textbf{Build système flexible} avec CMake supportant GPU, CPU, tests
\end{itemize}

Les perspectives d'amélioration incluent l'ajout de sources de lumière explicites,
le support des textures, l'implémentation du sampling par importance multiple (MIS),
et l'extension aux volumes participatifs (brouillard, fumée).

% -----------------------------------------------
\end{document}
